% !PRINTED 754
\weberpage{754}
Then $F(x)$ is an integral function of $x$ of degree $n$. If finally we also set

\weq{10}{Q(x) = c_n q_n x^n + c_{n-1}q_{n-1}x^{n-1}+\cdots+c_0 q_0,}

\weq{11}{P = c_n\Pi(n)+c_{n-1}\Pi(n-1)+\cdots+c_0,}

then $Q(x)$ depends on $x$, though it is not expressible rationally in terms of $x$; $P$, however, is independent of $x$.

If we now multiply equations (6) in succession by $c_n, c_{n-1},\ldots, c_0$ and add them, it follows that

\weq{12}{e^x P = F(x) + e^{\xi}Q(x).}

This formula is now the starting point of the further conclusions.

Let us assume that $e$ is an algebraic number; then an equation must subsist, whose degree let us call $m$:

\weq{13}{C_0 + C_1 e + C_2 e^2 + \cdots + C_m e^m = 0,}

whose coefficients $C_0, C_1, \ldots, C_m$ are \gesperrt{integral rational numbers}, among which $C_0$ and $C_m$ are different from zero. The matter is to demonstrate the impossibility of this assumption.

For this purpose we substitute in equation (12) for $x$ in succession the integers $0,1,2,\ldots,m$, so that $\xi$ becomes identical with $x$, multiply by $C_0, C_1, \ldots, C_m$ and add. Then, by (13), since $P$ is independent of $x$, there results:

\weq{14}{0 = C_0 F(0) + C_1 F(1) + \cdots + C_m F(m) + C_0 Q(0) + C_1 e\,Q(1) + \cdots + C_m e^m Q(m),}

and now it is to be shown, from a suitable assumption about the still arbitrary function $f(x)$, that equation (14) is impossible.

We choose a prime $p$ which is greater than $m$\,\footnote{That there is always a prime $p$ greater than an arbitrary number $\mu$ was already proved in \gesperrt{Euclid} (Elemente, Buch IX, Nr. XX, Bd. 2 of the \gesperrt{Heiberg} edition). The proof is simply this, that the integer $\Pi(\mu)+1$, which is evidently greater than $\mu$, is divisible by no prime that is not greater than $\mu$, because on division by each of the numbers $2,3,\ldots,\mu$ this number gives the residue 1. It is therefore impossible that no prime lies above $\mu$.}, and set

\weq{15}{f(x) = \frac{x^{p-1}(1-x)^p(2-x)^p\ldots(m-x)^p}{\Pi(p-1)},}
%CONT>>
